%% kapa: exam=1 round=32 year=2021 period=1 form=A subject=civil_law subject_ko=민법 seq=20 number=20
%% source: Q-Net 감정평가사 기출문제 (cst00309) artlSeq=5208397 file=extracted/1008/5208397/A형/2021년도 제32회 감정평가사 1차 1교시 A형.pdf
%% conversion: kordoc 4.24.51 -> tools/gen1.py (auto); review=proofread-sonnet; flags=none
\begin{question}{20}
소멸시효에 관한 설명으로 옳지 않은 것은? (다툼이 있으면 판례에 따름)
\begin{choices}
  \item 소멸시효는 법률행위에 의하여 이를 배제하거나 연장할 수 없다.
  \item 시효의 중단은 원칙적으로 당사자 및 그 승계인 사이에서만 효력이 있다.
  \item 소멸시효 중단사유로서의 채무승인은 채무가 있음을 알고 있다는 뜻의 의사표시이므로 효과의사가 필요하다.
  \item 소멸시효의 이익은 시효가 완성되기 전에 미리 포기하지 못한다.
  \item 소멸시효 완성 후 채무자는 시효완성의 사실을 알고 그 채무를 묵시적으로 승인함으로써 시효의 이익을 포기할 수 있다.
\end{choices}
\end{question}
\answer{3}
